\chapter{量子理论基础}
\section{Stern-Garlach实验}
\begin{center}
    \begin{minipage}{0.60\textwidth}
        \centering
        % 调整 x 比例让图形在 minipage 中更紧凑
        \begin{tikzpicture}[>=Stealth, line width=0.8pt]
            % -银原子源
            \node at (-3, 1.5) {\Large $47$ \textbf{Ag}};
            \foreach \pos in {
                (-3.4, 0.2), (-3.6, 0.5), (-3.4, 0.8), (-3.1, 0.6), (-2.8, 0.4), (-2.5, 0.6), (-2.3, 0.2),
                (-3.3, -0.1), (-3.5, -0.3), (-3.2, -0.4), (-2.9, -0.1), (-2.4, -0.2), (-2.2, -0.3),
                (-3.2, 0.4), (-2.9, 0.1), (-3.3, -0.2), (-2.6, -0.1), (-3.0, 0.5), (-2.7, 0.4)
            } {\fill \pos circle (2.2pt);}
            % 电子排布文字
            \node at (-3, -1) {\large $2, 8, 18, 18, 1$};
            \node at (-1.5, -2.2) {\large $5s: n=5, l=0, m=0$};
            % 准直孔
            \draw[line width=3pt] (-1, 1.5) -- (-1, 0.2);
            \draw[line width=3pt] (-1, -0.2) -- (-1, -1.5);
            % 入射束箭头
            \draw[->, OrangeDeep, line width=1.5pt] (-1.8, 0) -- (-0.2, 0);
            %使用 plot 绘制两条对称的分离曲线
            \begin{scope}[OrangeDeep, line width=1.2pt]
                \draw[domain=0:4, samples=50, smooth] plot ({0.7+\x}, {0.15*exp(0.6*\x) - 0.15});
                \draw[domain=0:4, samples=50, smooth] plot ({0.7+\x}, {-0.15*exp(0.6*\x) + 0.15});
                \draw (0.5, 0) -- (0.7, 0); % 磁场入口直线
            \end{scope}
            % 绘制向下递减的蓝色箭头
            \begin{scope}[shift={(1, 1.2)}, color=BlueDeep]
                \foreach \x in {0, 0.2, ..., 1.6} \draw[->] (\x, 0) -- (\x, -0.5);
                \foreach \x in {0.4, 0.6, ..., 1.4} \draw[->] (\x, -0.6) -- (\x, -1);
                \foreach \x in {0.6, 0.8, ..., 1.0} \draw[->] (\x, -1.1) -- (\x, -1.6);
                \draw[->] (0.8, -1.7) -- (0.8, -2.2);
                \node[below, font=\small] at (0.8, -2.2) {非均匀磁场};
            \end{scope}
            % 探测屏与坐标轴
            \draw[line width=3pt] (5, 1.8) -- (5, -1.8);
            \begin{scope}[shift={(3.8, 0)}]
                \draw[->] (0,0) -- (0.6, 0) node[right] {$x$};
                \draw[->] (0,0) -- (0, 0.6) node[above] {$z$};
                \draw[->] (0,0) -- (-0.25, -0.25) node[left] {$y$};
                \fill (0,0) circle (1.5pt); % 原点
            \end{scope}
        \end{tikzpicture}
    \end{minipage}
    \hfill
    \begin{minipage}{0.36\textwidth}
        \raggedright
        Ag轨道角动量为$0$，轨道磁矩为$0$.

        $z$方向受力$F_z=\mu_z \cdot \frac{\partial B}{\partial z}$，而银原子分裂成了两束.说明磁矩$\mu_z$不为$0$，且是量子化的.
    \end{minipage}
\end{center}
该实验表明电子除轨道角动量外还有一类角动量且它的大小为$\dfrac{\hbar}{2}$ (实验测定).

\subsection{电子自旋假设}
电子自旋假设由Uhlenbeck和Goudsmit于1925年提出：
\begin{itemize}
    \item 每个电子都具有内禀角动量(即\overtext{自旋}{spin}角动量)$\hat{S}$，大小为$\frac{\hbar}{2}$. \\
    为什么说是内禀的自旋？如果按传统的“旋转产生角动量”来理解，电子表明的切速度会超过光速.
    \item 自旋磁矩与自旋角动量有如下关系：
    \[
    \begin{cases}
        \hat{\mu}_s = -\frac{e}{2m_e}g_s \hat{S} \\
        g_s = 2
    \end{cases}
    \quad
    \left(
        \text{轨道磁矩}\begin{cases}
            \hat{\mu}_l = -\frac{e}{2m_e} g_l \hat{L} \\
            g_l = 1
        \end{cases}
    \right)
    \]
    向上偏转记为$\ket{\uparrow }$，反之为$\ket{\downarrow}$.
\end{itemize}
\begin{extension}{Bosen子和Fremi子}
    称整数自旋的粒子如光子$(1)$、希格斯子$(0)$为\textbf{Bosen子}.

    称非整数自旋的粒子如电子$(\frac{1}{2})$、夸克$(\frac{1}{2})$为\textbf{Fermi子}.
\end{extension}

\subsection{级联S-G实验}
\begin{center}
    \begin{tikzpicture}[
        >=Stealth, 
        line width=1.2pt,
        sgbox/.style={draw, minimum width=1.5cm, minimum height=1.5cm, rounded corners=2pt, align=center, font=\Large\bfseries}
    ]
        \node[sgbox] (SG1) at (0,0) {S-G \\ $z$};
        \draw[->] (SG1.east) ++(0,0.4) -- +(1.5,0) node[above, pos=0.45] {$\ket{\uparrow}$};
        \draw[->] (SG1.east) ++(0,-0.4) -- +(1.5,0) node[below, pos=0.45] {$\ket{\downarrow}$};

        \node[sgbox, anchor=west] (SG2) at (2.25, 0.8) {S-G \\ $x$};
        \draw[->] (SG2.east) ++(0,0.4) -- +(1.5,0) node[above, pos=0.45] {$\ket{+}$};
        \draw[->] (SG2.east) ++(0,-0.4) -- +(1.5,0) node[below, pos=0.45] {$\ket{-}$};

        \node[sgbox, anchor=west] (SG3) at (5.25, 1.6) {S-G \\ $z$};
        \draw[->] (SG3.east) ++(0,0.3) -- +(1.5,0) node[above, pos=0.45] {$\ket{\uparrow}$};
        \draw[->] (SG3.east) ++(0,-0.4) -- +(1.5,0) node[below, pos=0.45] {$\ket{\downarrow}$};
    \end{tikzpicture}
\end{center}
\noindent\textbf{实验现象：}
\begin{itemize}
    \item 银原子束通过第一个$z$方向的S-G后分为两束，取向上的一束通过$x$方向的S-G后分为两束，各占$50\%$.
    \item 取向+的一束通过第二个$z$方向的S-G，仍能分为向上和向下的两束，各占$50\%$.
\end{itemize}

这个实验表明：对于$S_x$的测量把粒子$S_z$的信息抹去了.

因此，无法同时得到关于粒子$S_x$和$S_z$的信息，称$S_x$与$S_z$是不相容的可观测量(由对称性，$S_x,S_y,S_z$都不相容).
\begin{align*}
    &\text{$\ket{\uparrow}$里包含了$\ket{+}$和$\ket{-}$的成分} \quad &\ket{\uparrow} = \frac{1}{\sqrt{2}} e^{i\theta} \ket{+} + \frac{1}{\sqrt{2}} e^{i\theta} \ket{-} \\
    &\text{$\ket{+}$里包含了$\ket{\uparrow}$和$\ket{\downarrow}$的成分} &\ket{+} = \frac{1}{\sqrt{2}} e^{i\theta} \ket{\uparrow} + \frac{1}{\sqrt{2}} e^{i\theta} \ket{\downarrow}
\end{align*}
\begin{mistake}
    上述表述似乎意味着$\ket{\uparrow}$和$\ket{+}$既是本征态又是叠加态.

    实际上，讨论一个态是什么态要\textbf{基于某个对象}才有意义：
    \begin{align*}
        \text{对$S_x$，$\ket{\uparrow}$是叠加态，$\ket{+}$是本征态} \\
        \text{对$S_z$，$\ket{\uparrow}$是本征态，$\ket{+}$是叠加态}
    \end{align*}

    这表明：

    \centering\textbf{在经典和量子世界中，测量扮演着十分不同的角色}
\end{mistake}


\section{基于线性代数的量子力学I}
\subsection{量子力学与线性代数的相似性}
\begin{enumerate}
    \item 可观测量在其本征态上有一个确定的值. \\
    矩阵作用在本征矢上有一个确定的值.
    \item 不同的态可以线性叠加成为一个新的可能态. \\
    矢量之间可以通过线性叠加构造一个新的矢量.
    \item Schrödinger方程是线性偏微分方程，其通解为本征解(特解)的线性组合.
    $$\psi(\vec{r}, t) = \sum_{i} c_i \phi_i(\vec{r}) T_i(t)$$
    这些本征解可以作为LVS的基，任意的通解可以通过$c_i$的线性组合来表述：
    $$\lvert \psi \rangle = \sum_{i} c_i T_i(t) \lvert \phi_i \rangle = \sum_{i} c_i \lvert \phi_i(t) \rangle = \begin{bmatrix} c_1 \\ c_2 \\ \vdots \\ c_n \end{bmatrix}$$
    \item 可观测量由初、末两个态决定，应是二维的，而测量会使一个态转化为另一个态.\\
    矩阵是二阶张量，矩阵作用在一个矢量上会使之变为另一个矢量.
\end{enumerate}

综上，用矢量来描述量子系统的状态，即态矢量$\ket{\psi}$，它包含了系统的全部信息；用矩阵来描述力学量$\hat{O}$，力学量是物理的，其测量值(本征值)应为实数，因此描述力学量的矩阵应为厄米矩阵.

\subsection{拓展到无穷维}
\begin{question}
    考虑定义在$[a,b]$上的函数$f(x)$，如何用一个列矢量描述$f(x)$？

    \begin{center}
        \begin{tikzpicture}[>=stealth, scale=0.8]
            \draw[->] (-0.5,0) -- (6,0) node[right] {$x$};
            \draw[->] (0,-2) -- (0,2) node[right] {$f(x)$};
            \draw[OrangeDeep, thick, samples=100, domain=1:5.5] 
                plot (\x, {1.5*sin((\x-1)*120) * (1.2 - 0.15*\x)});
            \node[below] at (1,0) {$a$};
            \node[below] at (5.5,0) {$b$};
        \end{tikzpicture}
    \end{center}
\end{question}
\begin{solution}
    将$a~b$分成$n$份(生成$n+1$个基矢)，用每处$x_i$的函数值$f(x_i)$作列矢量的矩阵元：
    $$|f\rangle = \begin{bmatrix} f(x_1) \\ f(x_2) \\ \vdots \\ f(x_n) \end{bmatrix}$$
    $f(x)$在每一处$x_i$的函数值等于$\ket{f}$沿着$\ket{x_i}$的分量，即$f(x_i) = \braket{x_i}{f}$.
\end{solution}

这意味着：
\begin{center}
    \textbf{波函数是态矢量在具体基(表象)下的展开系数}
\end{center}
考虑以$\ket{x_i}$为基张成的空间满足LVS的全部性质，那么$\ket{x_i}$的所有线性组合构成的矢量构成一个LVS，维度为$n+1$.当$n \rightarrow \infty$时构成了一个连续无穷维LVS.

\subsection{构造连续无穷维LVS的IPS}
在离散有限维LVS中，内积的定义为：
$$\langle f | g \rangle = \lim_{n \to \infty} \sum^{n}_{i} f^*(x_i)g(x_i)$$
由于$f(x_i)$和$g(x_i)$都是有限的，因此在无穷维上求和会是无穷大.这导致在连续无穷维LVS上如矩阵乘法等数学工具的发散，因此，需要构建一套新的运算规则.
\begin{derivation}
    为了使得结果有界，取内积在每个点上的平均：
    \begin{align*}
        \text{令} \quad \langle f|g \rangle &= \lim_{n \to \infty} \frac{1}{n+1} \sum_{i} f^*(x_i) g(x_i) \\
        &= \int_a^b f^*(x) g(x) \dd{x} \\
        &= \int_a^b \langle f|x \rangle \langle x|g \rangle \dd{x}
    \end{align*}
    为不失一般性，选取积分区域为整个空间：
    $$\langle f|g \rangle = \int_{-\infty}^{\infty} \langle f|x \rangle \langle x|g \rangle \dd{x}$$
    这就是连续无穷维LVS内积的定义.

    称式中的$\int_{-\infty}^{\infty} |x \rangle \langle x| \dd{x} \equiv \hat{I}$为连续无穷维情况下的恒等算符.
\end{derivation}

\subsection{Hilbert空间}
量子体系是物理的且可能是无穷维的.

态矢量应当满足如下需求：

\begin{itemize}
    \item 态矢量可以进行内积
    \item 态矢量可以是有限维的也可以是无穷维的
    \item 态矢量之间的内积有限
\end{itemize}

因此量子力学所研究的空间应当满足：

\begin{itemize}
    \item 是一个IPS
    \item 维度可以从有限到无限
    \item 其上定义的内积是收敛的
\end{itemize}

称这样的空间为\textbf{Hilbert空间}.

\begin{extension}{}
    \textbf{应当注意：}
    \begin{itemize}
        \item 对同一个LVS，通过定义不同的内积法则可以构造出不同的IPS.
        \item Hilbert空间的严格定义是“完备的内积空间”.更深入的内容涉及泛函分析.
        \item 量子力学中会有一些内积为奇怪东西(如Dirac $\delta$函数)的量子态，这显然不存在于数学意义上的Hilbert空间中，但是在物理上是可以接受的.称这种包含奇怪量子态的空间为rigged Hilbert空间.
    \end{itemize}
\end{extension}


\section{Dirac \texorpdfstring{$\delta$}{delta} 函数}
\subsection{Dirac归一化}
考察连续无穷维LVS中基矢的正交归一性：
\begin{derivation}
    作为基矢，显然有$\braket{x}{x'}=0, x \neq x'$.

    由于定义了新的IPS，$\braket{x}{x'}$在$x=x'$时是否等于$1$？

    将$\ket{f}$用$\ket{x'}$展开：
    \begin{align*}
        |f\rangle &= \int_{a}^{b} |x'\rangle \langle x'|f\rangle \dd{x'}
        \intertext{同时左乘$\bra{x}$：}
        \langle x|f\rangle &= \int_{a}^{b} \langle x|x'\rangle \langle x'|f\rangle \dd{x'} \\
        f(x) &= \int_{a}^{b} \langle x|x'\rangle f(x') \dd{x'} \\
        \intertext{由于有$\langle x|x'\rangle = 0, \quad x \neq x'$，积分区间可写为$x$附近的无穷小邻域：}
        f(x) &= \int_{x-\varepsilon}^{x+\varepsilon} \langle x|x'\rangle f(x') \dd{x'} \\
        \intertext{在无穷小区间上$f(x)$与$f(x')$近似相等：}
        f(x) &= \int_{x-\varepsilon}^{x+\varepsilon} \langle x|x'\rangle f(x) \dd{x'} \\
        \intertext{$f(x)$与关于$x'$的积分无关，作常数提出：}
        f(x) &= f(x) \int_{x-\varepsilon}^{x+\varepsilon} \langle x|x'\rangle \dd{x'} \quad \Rightarrow \quad \int_{x-\varepsilon}^{x+\varepsilon} \langle x|x'\rangle \dd{x'} = 1
    \end{align*}
    要在无穷小区间上积分为有限值，则$\braket{x}{x'}$在积分区间上不能是有限值，则:
    $$\langle x|x'\rangle = \infty, \quad x-\varepsilon \le x \le x+\varepsilon$$
    $$\begin{cases} \langle x|x'\rangle = 0 & , x' \neq x \\ \langle x|x'\rangle \to \infty & , x' = x \end{cases}$$
\end{derivation}
\begin{definition}
    Dirac $\delta$函数$\braket{x}{x'}= \delta (x-x')$满足：
    \[
    \delta(x-x')=
    \begin{cases}
        0, x\neq x'\\
        \infty, x\rightarrow x'
    \end{cases}
    \qquad
    \int_{-\infty}^{+\infty} \delta(x-x') \dd{x} = 1
    \]
    图像表示为：
    \begin{center}
        \begin{tikzpicture}[>=stealth, font=\small]
            \begin{scope}[shift={(0,0)}]
                \draw[->, thick] (0,-0.2) -- (0,3.5);
                \draw[->, thick] (-0.2,0) -- (4,0) node[below] {$x'$};
                \draw[YellowDeep, very thick] plot[domain=0.2:3.8, samples=100] (\x, {0.004 / ((\x-2)^2 + 0.001)});
                \node[YellowDeep, below] at (2,0) {$x$};
            \end{scope}
            \node at (5, 1.5) {\Large 或};
            \begin{scope}[shift={(6,0)}]
                \draw[->, thick] (0,-0.2) -- (0,3.5);
                \draw[->, thick] (-0.2,0) -- (4,0) node[below] {$x'$};
                \draw[YellowDeep, ultra thick, ->] (2,0) -- (2,3);
                \node[YellowDeep, below] at (2,0) {$x$};
            \end{scope}
        \end{tikzpicture}
    \end{center}
    显然$\delta(x)$是关于$x$的偶函数.
\end{definition}

\subsection{Dirac \texorpdfstring{$\delta$}{delta} 函数的性质}
\begin{itemize}
    \item \textbf{采样性}\\
    把一个函数与$\delta(x-x')$相乘并在全空间积分，就可以采集到该函数在$x_0$处的值：
    $$\int_{-\infty}^{+\infty} f(x) \delta(x-x_0) \dd{x} = f(x_0)$$
    \item \textbf{放缩性}
    $$\delta(ax) = \frac{1}{|a|} \delta(x)$$
    \begin{proof}
        \begin{minipage}[t]{0.48\textwidth}
            \centering
            \textbf{当$a > 0$时}
            \begin{align*}
                & \int_{-\infty}^{+\infty} f(x) \delta(ax) \dd x \quad \text{令 } t \equiv ax \\
                = & \int_{-\infty}^{+\infty} f\left(\frac{t}{a}\right) \delta(t) \dd \left(\frac{t}{a}\right) \\
                = & \frac{1}{a} \int_{-\infty}^{+\infty} f\left(\frac{t}{a}\right) \delta(t) \dd t \\
                \intertext{令$F(t)\equiv f(\frac{t}{a})$，有$F(0)=f(0)$} 
                = & \frac{1}{a} F(0) = \frac{1}{a} f(0)
            \end{align*}
        \end{minipage}
        \hfill\vrule\hfill
        \begin{minipage}[t]{0.48\textwidth}
            \centering
            \textbf{当$-a > 0$时}
            \begin{align*}
                & \int_{-\infty}^{+\infty} f(x) \delta(ax) \dd x \\
                = & \int_{+\infty}^{-\infty} f\left(-\frac{t}{a}\right) \delta(t) \dd \left(-\frac{t}{a}\right) \\
                = & -\frac{1}{a} \int_{+\infty}^{-\infty} f\left(-\frac{t}{a}\right) \delta(t) \dd t \\
                = & \frac{1}{a} \int_{-\infty}^{+\infty} F(-t) \delta(t) \dd t \\
                = & \frac{1}{a} F(0) = \frac{1}{a} f(0)
            \end{align*}
        \end{minipage}\\[1em]
    综上，$\delta(ax) = \frac{1}{|a|} \delta(x)$.
    \end{proof}
    \item \textbf{实值性}
    $$\langle x | x' \rangle = \delta(x-x') = \delta(x'-x) = \langle x' | x \rangle = \langle x | x' \rangle^*$$
    因此Dirac $\delta$函数为实值函数.
\end{itemize}

\subsection{Dirac \texorpdfstring{$\delta$}{delta} 函数的Fourier变换}
Fourier变换的标准形式：
$$\begin{cases} 
\text{正变换 } T(k) = \frac{1}{\sqrt{2\pi}} \int_{-\infty}^{+\infty} e^{-ikx} t(x) \dd{x} \\
\text{逆变换 } t(x) = \frac{1}{\sqrt{2\pi}} \int_{-\infty}^{+\infty} e^{ikx} T(k) \dd{k}
\end{cases}$$
由于：
$$\frac{1}{\sqrt{2\pi}} \int_{-\infty}^{+\infty} e^{-ikx'} \delta(x - x') \dd{x'} = \frac{1}{\sqrt{2\pi}} e^{-ikx}$$
故而：
\begin{align*}
    \delta(x - x') &= \frac{1}{\sqrt{2\pi}} \int_{-\infty}^{+\infty} e^{ikx'} \cdot \frac{1}{\sqrt{2\pi}} e^{-ikx} \dd{k} \\
    \delta(x - x') &= \frac{1}{2\pi} \int_{-\infty}^{+\infty} e^{ik(x' - x)} \dd{k} \quad \delta \text{(函数的 Fourier 展开)}
\end{align*}

\subsection{含Dirac \texorpdfstring{$\delta$}{delta} 函数的导数的积分}
\begin{derivation}
    $$\begin{aligned}
        & \int_{-\infty}^{+\infty} \delta''(x - x') f(x') \dd{x'} \\
        = & \int_{-\infty}^{+\infty} \frac{\partial^2}{\partial x^2} \delta(x - x') f(x') \dd{x'} \\
        = & \frac{\dd^2}{\dd{x^2}} \int_{-\infty}^{+\infty} \delta(x - x') f(x') \dd{x'} \quad \text{积分外} x' \text{被积掉了，故而变为求全导} \\
        = & \frac{\dd^2}{\dd{x^2}} f(x)
    \end{aligned}$$
\end{derivation}

上述推导表明，$\delta$函数的采样性亦可给出$f(x)$的各阶导：
\begin{gather*}
\delta^{(n)}(x - x') = \delta(x - x') \frac{\dd^n}{\dd{x^n}} \\
\int_{-\infty}^{+\infty} \delta^{(n)}(x - x') f(x') \dd{x'} = \frac{\dd^n}{\dd{x^n}} f(x)
\end{gather*}

\subsection{单位跃阶函数}
\begin{definition}
    \[
    \text{单位跃阶函数}\theta(x)=
    \begin{cases}
        1,x>0 \\
        0,x<0
    \end{cases}
    \qquad
    \frac{\dd}{\dd{x}}\theta(x)=\delta(x)
    \begin{cases}
        \infty,x=0 \\
        0,x\neq 0
    \end{cases}
    \]
    \begin{center}
        \begin{tikzpicture}[>=stealth, thick]
            \draw[->] (-3,0) -- (3,0) node[below] {$x$};
            \draw[->] (0,-0.5) -- (0,2) node[right] {$\theta$};
            \node[below left] at (0,0) {0};
            \draw (0.1,1) -- (-0.1,1) node[left] {1};
            \draw[YellowDeep, ultra thick] (-2.5,0) -- (0,0);
            \draw[YellowDeep, ultra thick] (0,1) -- (2.5,1);
        \end{tikzpicture}
    \end{center}
    关于$x=0$处的定义不同领域各不相同，物理学上一般定义为$\frac{1}{2}$.
\end{definition}


\section{基于线性代数的量子力学II}
\subsection{无穷维下的算符}
无穷维下的右矢对应了一个连续函数，因此：
\begin{center}
    \textbf{无穷维的算符就是将一个函数$f(x)$转换为另一个函数$\widetilde{f}(x)$}
\end{center}
即$\hat{O}\ket{f}=\ket{\widetilde{f}}$.

\begin{extension}{求导算符$\hat{D}$}
    对求导算符有：
    \[\hat{D}\ket{f}=\ket{\frac{\dd}{\dd x} f} \qquad \braket{x}{\frac{\dd}{\dd x}f}=\frac{\dd f}{\dd x}\]
    要求$\hat{D}$在基$\{\ket{x}\}$($x$表象)下的矩阵表示，就是要求$\hat{D}$的矩阵元$D_{xx'} = \bra{x}\hat{D}\ket{x'}$.
    \begin{align*}
        \hat{D}\ket{f} &= \ket{\frac{\dd}{\dd x}f} \\
        \bra{x}\hat{D}\ket{f} &= \bra{x}\ket{\frac{\dd}{\dd x} f} \\
        \int_{-\infty}^{+\infty}\bra{x}\hat{D}\ket{x'}\bra{x'}\ket{f} \dd{x'} &= \bra{x}\ket{\frac{\dd}{\dd{x'}}} = \ket{x}\ket{\frac{\dd}{\dd x} f} \\
        \int_{-\infty}^{+\infty} D_{xx'} f(x') \dd{x'} &= f'(x) \\
        \Rightarrow \delta'(x-x') = \delta(x-x')\frac{\dd}{\dd{x'}}
    \end{align*}
\end{extension}
\begin{proposition}
    求导算符是反厄米算符.
\end{proposition}
\begin{proof}
    即证明$\hat{D}^{\dagger} = (\hat{D}^{*})^{\top} = -\hat{D} \Rightarrow D_{xx'} = -D^*_{x'x}$.
    \begin{align*}
        D^*_{x'x} &=  [\delta ' (x'-x)]^* \quad \text{$\delta$函数的实值性}\\
        &= \delta ' (x'-x) \quad \text{偶函数的一阶导是奇函数} \\
        &= -\delta ' (x-x') \\
        &= -D_{xx'}
    \end{align*}
\end{proof}

可见，由反厄米算符构造厄米算符只需乘以一个$i$.

\begin{extension}{位置算符$\hat{X}$}
    有本征方程：
    \begin{align*}
        \hat{X}\ket{x_j} &= x_j \ket{x_j} \\
        \bra{a_i}\hat{X}\ket{x_j} &= x_j\braket{x_i}{x_j} \\
        X_{ij} &= x_j \delta(x_i - x_j)
    \end{align*}
    因此$X$在$x$表象下的矩阵为对角阵.
\end{extension}

\subsection{无穷维下的厄米算符}
\begin{proposition}
    若$\hat{K}$为厄米算符，则$\braket{\hat{K}f}{g} = \bra{f}\hat{K}\ket{g} = \braket{f}{\hat{K}g}$.
\end{proposition}
\begin{proof}
    $$\langle \hat{K} f | = ( | \hat{K} f \rangle )^{\dagger} = \langle f | \hat{K} \implies \langle \hat{K} f | g \rangle = \langle f | \hat{K} g \rangle$$
\end{proof}
这是厄米算符的另一种判据，在不易求得$\hat{K}^{\top}$时可用.
\begin{question}
    说明$\hat{K} = -i \frac{\dd}{\dd x}$在什么情况下是厄米算符.
\end{question}
\begin{solution}
    根据厄米算符的性质：
    $$\langle \hat{K} f | g \rangle = \langle f | \hat{K} g \rangle$$
    \begin{align*}
        \text{左边} &= \int_a^b (\hat{K} f)^* g \dd x \\
        &= \int_a^b i \frac{\dd}{\dd x} f^* g \dd x \\
        &= \int_a^b i g \dd{f^*} \\
        &= i g f^* \big|_a^b - \int_a^b i f^* \dd g \\
        &= i g f^* \big|_a^b + \int_a^b f^* (-i \frac{\dd}{\dd x} g) \dd x \\
        &= \langle f | \hat{K} g \rangle + i g f^* \big|_a^b
    \end{align*}
    显然，只有边界项$i f^* g \big|_a^b = 0$时$\hat{K} = -i \frac{\dd}{\dd x}$才是厄米算符.
\end{solution}

这表明一个算符是否为厄米算符还与其所作用的态有关.

\subsection{力学量的平均值}
以位置为例，$|\psi(x)|^2$表示粒子在$x$处出现的概率密度$\rho(x)$.

$x$在全空间上的均值：
\begin{align*}
    \bar{x} &= \int_{-\infty}^{+\infty} x \rho(x) \dd x \\
    &= \int_{-\infty}^{+\infty} x |\psi(x)|^2 \dd x \\
    &= \int_{-\infty}^{+\infty} x \cdot \psi^(x) \cdot \psi(x) \dd x \\
    &= \int_{-\infty}^{+\infty} \psi^(x) \cdot x \cdot \psi(x) \dd x \\
    &= \langle \psi | x | \psi \rangle \quad \text{记作 } \langle x \rangle
\end{align*}

对任意力学量$\hat{\Omega}$，在$\ket{\psi}$态下对其测量时总是会落入对应特定本征值$\omega_i$的本征态$\ket{\omega_i}$，相应的概率为$p_i = |\langle \omega_i | \psi \rangle|^2$.

力学量$\hat{\Omega}$在某个态$\ket{\psi}$下的平均值为：
\begin{align*}
    \bar{\Omega} &= \sum_i p_i \omega_i \\
    &= \sum_i |\langle \omega_i | \psi \rangle|^2 \omega_i \\
    &= \sum_i \langle \psi | \omega_i \rangle \langle \omega_i | \psi \rangle \omega_i \\
    &= \sum_i \langle \psi | \omega_i | \omega_i \rangle \langle \omega_i | \psi \rangle \quad (\hat{\Omega} \text{ 的本征方程 } \hat{\Omega} |\omega_i\rangle = \omega_i |\omega_i\rangle) \\
    &= \sum_i \langle \psi | \hat{\Omega} | \omega_i \rangle \langle \omega_i | \psi \rangle \\
    &= \langle \psi | \hat{\Omega} | \psi \rangle \quad \text{记作 } \langle \hat{\Omega} \rangle
\end{align*}

\subsection{对应原理}
Bohr认为在微观、少粒子的情况下，体系表现出量子力学特性；当逐渐趋于宏观，大量粒子的情况时。其统计结果逐步趋向于经典力学.
\begin{derivation}
    自由粒子速度的平均值：
    \begin{align*}
        \bar{v} &= \frac{\dd \bar{x}}{\dd t} \\
        &= \frac{\dd}{\dd t} \int_{-\infty}^{+\infty} x |\psi(x,t)|^2 \dd x \\
        &= \int_{-\infty}^{+\infty} \left( \frac{\partial}{\partial t} x \cdot |\psi(x,t)|^2 + x \cdot \frac{\partial}{\partial t} |\psi(x,t)|^2 \right) \dd x \\
        &= \int_{-\infty}^{+\infty} x \cdot \frac{\partial}{\partial t} |\psi(x,t)|^2 \dd x \\
        &= \int_{-\infty}^{+\infty} x \cdot \underbrace{\left[ \psi \frac{\partial}{\partial t} \psi^* + \psi^* \frac{\partial}{\partial t} \psi \right]}_{\text{代入Schrödinger方程}} \dd x \\
        &= \int_{-\infty}^{+\infty} x \cdot \left[ \psi \left( \frac{i}{\hbar} \hat{H} \psi^* \right) - \psi^* \left( \frac{i}{\hbar} \hat{H} \psi \right) \right] \dd x \\
        &= \int_{-\infty^{+\infty}} x \cdot \frac{i}{\hbar} \cdot \frac{\hbar^2}{2\mu} \left[ \psi^* \frac{\partial^2}{\partial x^2} \psi - \frac{\partial^2}{\partial x^2} \psi^* \cdot \psi \right] \dd x \\
        &= \int_{-\infty}^{+\infty} x \cdot \frac{i\hbar}{2\mu} \cdot {\frac{\partial}{\textcolor{BlueDeep}{\partial x}} \left[ \psi^* \frac{\partial}{\partial x} \psi - \frac{\partial}{\partial x} \psi^* \cdot \psi \right] \textcolor{BlueDeep}{\dd x}} \\
        &= \int_{-\infty}^{+\infty} x \cdot \frac{i\hbar}{2\mu} \textcolor{BlueDeep}{\dd \left[ \psi^* \frac{\partial}{\partial x} \psi - \frac{\partial}{\partial x} \psi^* \cdot \psi \right]} \\
        &= \frac{i\hbar}{2\mu} \left\{\textcolor{OrangeDeep}{\underbrace{\left[ x \cdot \left( \psi^* \frac{\partial}{\partial x} \psi - \frac{\partial}{\partial x} \psi^* \cdot \psi \right) \right]_{-\infty}^{+\infty}}_{\text{波函数在无穷远处为$0$}}} - \int_{-\infty}^{+\infty} \left[ \psi^* \frac{\partial}{\partial x} \psi - {\frac{\partial}{\textcolor{BlueDeep}{\partial x}} \psi^* \cdot \psi} \right] \textcolor{BlueDeep}{\dd x} \right\} \\
        &= \textcolor{OrangeDeep}{0} - \frac{i\hbar}{2\mu} \left\{ \int_{-\infty}^{+\infty} \left[ \psi^* \frac{\partial}{\partial x} \psi \right] \dd x - \int_{-\infty}^{+\infty} \psi \dd \psi^* \right\} \\
        &= -\frac{i\hbar}{2\mu} \left\{ \int_{-\infty}^{+\infty} \left[ \psi^* \frac{\partial}{\partial x} \psi \right] \dd x - \left[ \textcolor{OrangeDeep}{\psi^* \psi \Big|_{-\infty}^{+\infty}} - \textcolor{PurpleDeep}{\int_{-\infty}^{+\infty} \psi^* \dd \psi} \right] \right\} \\
        &= -\frac{i\hbar}{2\mu} \left\{ \int_{-\infty}^{+\infty} \left[ \psi^* \frac{\partial}{\partial x} \psi \right] \dd x + \textcolor{PurpleDeep}{\int_{-\infty}^{+\infty} \left[ \psi^* \frac{\partial}{\partial x} \psi \right] \dd x} \right\} \\
        &= \int_{-\infty}^{+\infty} \psi^* \frac{-i\hbar}{\mu} \frac{\partial}{\partial x} \psi \dd x \\
        &= \left\langle \psi \middle| \frac{-i\hbar}{\mu} \frac{\partial}{\partial x} \middle| \psi \right\rangle \\
        &= \langle v_x \rangle
    \end{align*}
    令：
    $$\hat{v}_x \equiv -\frac{i\hbar}{\mu}\frac{\partial}{\partial x}$$
    则动量算符：
    $$\hat{p} = -i\hbar \frac{\partial}{\partial x}$$
    它是一个厄米算符.
\end{derivation}

\subsection{量子力学的五大假设}
\begin{enumerate}
    \item \textbf{态假设} \\
    粒子的状态由Hilbert空间中的态矢量$\ket{\psi(t)}$完备地描述.
    \item \textbf{力学量假设} \\
    力学量由厄米算符描述.
    \item \textbf{测量假设} \\
    在$\ket{\psi}$态中对力学量$\hat{\Omega}$进行测量，会得到其本征值$\omega_i$，归一化后相应的概率为$|\braket{\omega_i}{\psi}|^2$，其中$\ket{\omega_i}$为$\omega_i$所对应的本征态.
    \item \textbf{演化假设} \\
    态的演化服从含时Schrödinger方程$i\hbar \frac{\partial}{\partial t}\ket{\psi(t)} = \hat{H}\ket{\psi(t)}$.
    \item \textbf{全同粒子假设} \\
    内禀属性完全相同的粒子不可分辨.
\end{enumerate}


\section{自旋1/2系统的矩阵表示}
在$\hat{S}_z$的本征矢所构成的基下，$\hat{S}_z$是一个对角阵.取：
$$\ket{\uparrow} = \begin{bmatrix} 1 \\ 0 \end{bmatrix},\qquad \ket{\downarrow} = \begin{bmatrix} 0 \\ 1 \end{bmatrix}$$
则在$\hat{S}_z$表象下：
$$\hat{S}_z = \begin{bmatrix} \frac{\hbar}{2} & 0 \\ 0 & -\frac{\hbar}{2} \end{bmatrix} = \underbrace{\frac{\hbar}{2} \ket{\uparrow}\bra{\uparrow} + (-\frac{\hbar}{2}) \ket{\downarrow}\bra{\downarrow}}_{\text{这样的表达不依赖于表象，是抽象的}}$$
而在$\hat{S}_x$表象下：
$$\hat{S}_x = \begin{bmatrix} \frac{\hbar}{2} & 0 \\ 0 & -\frac{\hbar}{2} \end{bmatrix} = \frac{\hbar}{2} \ket{+}\bra{+} + (-\frac{\hbar}{2}) \ket{-}\bra{-}$$

在表象下对力学量进行描述，只是用所属该表象的数学语言描述客观实在，而力学量作为客观实在是不依赖于表象而存在的.

\subsection{在\texorpdfstring{$\hat{S}_z$}{Sz}表象中表达\texorpdfstring{$\hat{S}_x$}{Sx}的矩阵}
\textbf{整体相位因子没有物理意义}，对波函数，乘以整体相位因子$e^{i\alpha}$后对其归一性没有任何影响：
\begin{align*}
    \rho = |\psi(x)|^2 &= \psi^*(x)\psi(x) \\
    &= (\psi(x)e^{i\alpha})^* \psi(x)e^{i\alpha} \\\
    &= |\psi(x)e^{i\alpha}|^2
\end{align*}
这样的特性也适用于态矢量，因为波函数的本质是态矢量在坐标基上的投影.
\begin{derivation}
    在级联S-G实验中有叠加态：
    $$\ket{\uparrow} = \frac{1}{\sqrt{2}} e^{i\theta} \ket{+} + \frac{1}{\sqrt{2}} e^{i\varphi} \ket{-}$$
    乘以一个整体相位因子$e^{i\theta}$，物理意义不变.
    $$\ket{\uparrow} = \frac{1}{\sqrt{2}} \ket{+} + \frac{1}{\sqrt{2}} e^{i(\varphi-\theta)} \ket{-}$$
    令$\delta_1 \equiv \varphi  - \theta$，则上式写为：
    $$\ket{\uparrow} = \frac{1}{\sqrt{2}} \ket{+} + \frac{1}{\sqrt{2}} e^{i\delta_1} \ket{-}$$
    同理有：
    $$\begin{cases} \ket{+} = \frac{1}{\sqrt{2}} \ket{\uparrow} + \frac{1}{\sqrt{2}} e^{i\delta_2} \ket{\downarrow} \\ \ket{-} = \frac{1}{\sqrt{2}} \ket{\uparrow} + \frac{1}{\sqrt{2}} e^{i\delta_2'} \ket{\downarrow} \end{cases}$$
    厄米算符对不同本征值的本征矢是相互正交的，因此$\braket{-}{+}=0$，即：
    \begin{align*}
        \braket{-}{+} &= \left( \frac{1}{\sqrt{2}} \ket{\uparrow} + \frac{1}{\sqrt{2}} e^{i\delta_2'} \ket{\downarrow} \right)^{\dagger} \cdot \left( \frac{1}{\sqrt{2}} \ket{\uparrow} + \frac{1}{\sqrt{2}} e^{i\delta_2} \ket{\downarrow} \right) \\
        &= \frac{1}{2} + \frac{1}{2} e^{i(\delta_2 - \delta_2')} \\
        &= 0
    \end{align*}
    则$e^{i\delta_2'} = -e^{i\delta_2}$.
    
    转换为相同的相对相位因子是为了确保它们“方向一致”，减少无贡献物理量.

    于是有:
    $$\begin{cases} \ket{+} = \frac{1}{\sqrt{2}} \ket{\uparrow} + \frac{1}{\sqrt{2}} e^{i\delta_1} \ket{\downarrow} \\ \ket{-} = \frac{1}{\sqrt{2}} \ket{\uparrow} - \frac{1}{\sqrt{2}} e^{i\delta_1} \ket{\downarrow} \end{cases}$$
    称之为$\hat{S}_x$的本征矢在$\hat{S}_z$表象下的形式.

    在$\hat{S}_x$表象下：
    $\hat{S}_x = \frac{\hbar}{2} \ketbra{+}{+} + (-\frac{\hbar}{2}) \ketbra{-}{-}$
    为在$\hat{S}_z$表象下表达$\hat{S}_x$，即须求$\ketbra{+}{+}$和$\ketbra{-}{-}$基于$\ket{\downarrow}$和$\ket{\uparrow}$的表达：
    \begin{align*}
        \ketbra{+}{+} =& \left( \frac{1}{\sqrt{2}} \ket{\uparrow} + \frac{1}{\sqrt{2}} e^{i\delta_1} \ket{\downarrow} \right) \cdot \left( \frac{1}{\sqrt{2}} \ket{\uparrow} + \frac{1}{\sqrt{2}} e^{-i\delta_1} \ket{\downarrow} \right)^{\dagger} \\
        =& \frac{1}{2} \left( \ketbra{\uparrow}{\uparrow} + e^{-i\delta_1} \ketbra{\uparrow}{\downarrow} + e^{i\delta_1} \ketbra{\downarrow}{\uparrow} + \ketbra{\downarrow}{\downarrow} \right) \\
        \ketbra{-}{-} =& \left( \frac{1}{\sqrt{2}} \ket{\uparrow} - \frac{1}{\sqrt{2}} e^{i\delta_1} \ket{\downarrow} \right) \cdot \left( \frac{1}{\sqrt{2}} \ket{\uparrow} - \frac{1}{\sqrt{2}} e^{-i\delta_1} \ket{\downarrow} \right)^{\dagger} \\
        =& \frac{1}{2} \left( \ketbra{\uparrow}{\uparrow} - e^{-i\delta_1} \ketbra{\uparrow}{\downarrow} - e^{i\delta_1} \ketbra{\downarrow}{\uparrow} + \ketbra{\downarrow}{\downarrow} \right) \\
        \hat{S}_x =& \frac{\hbar}{2} \cdot \frac{1}{2} \left( \ketbra{\uparrow}{\uparrow} + e^{-i\delta_1} \ketbra{\uparrow}{\downarrow} + e^{i\delta_1} \ketbra{\downarrow}{\uparrow} + \ketbra{\downarrow}{\downarrow} \right) \\
        &- \frac{\hbar}{2} \cdot \frac{1}{2} \left( \ketbra{\uparrow}{\uparrow} - e^{-i\delta_1} \ketbra{\uparrow}{\downarrow} - e^{i\delta_1} \ketbra{\downarrow}{\uparrow} + \ketbra{\downarrow}{\downarrow} \right) \\
        =& \frac{\hbar}{2} \left( e^{-i\delta_1} \ketbra{\uparrow}{\downarrow} + e^{i\delta_1} \ketbra{\downarrow}{\uparrow} \right) \\
        =& \frac{\hbar}{2} \begin{bmatrix} 0 & e^{-i\delta_1} \\ e^{i\delta_1} & 0 \end{bmatrix}
    \end{align*}
    接下来求取$\delta_2$的值，先考虑$\hat{S}_y$在$\hat{S}_z$表象下的形式.
    
    由于$y$轴和$x$轴在物理地位上是平权的，所以$\hat{S}_y$和$\hat{S}_z$在$\hat{S}_x$表象下的数学形式应为一致的：
    $$\hat{S}_y = \frac{\hbar}{2} \begin{bmatrix} 0 & e^{-i\delta_3} \\ e^{i\delta_3} & 0 \end{bmatrix}$$
    由级联S-G实验有：
    \begin{align*}
        |\braket{\textcolor{OrangeDeep}{S_{y+}}}{\textcolor{BlueDeep}{{-}}}|^2 &= \frac{1}{2} \\
        \left| \textcolor{OrangeDeep}{\left( \frac{1}{\sqrt{2}} \bra{\uparrow} + \frac{1}{\sqrt{2}} e^{-i\delta_3} \bra{\downarrow} \right)} \textcolor{BlueDeep}{\left( \frac{1}{\sqrt{2}} \ket{\uparrow} - \frac{1}{\sqrt{2}} e^{i\delta_2} \ket{\downarrow} \right)} \right|^2 &= \frac{1}{2} \\
        \frac{1}{4} |1 - e^{i(\delta_2 - \delta_3)}|^2 &= \frac{1}{2} \\
        |1 - e^{i(\delta_2 - \delta_3)}|^2 &= 2
    \end{align*}
    又由$\braket{S_{y-}}{-} = \frac{1}{2}$有$|1 + e^{i(\delta_2 - \delta_3)}|^2 = 2$，因此：
    $$\delta_2 - \delta_3 = \pm \frac{\pi}{2}$$
    \[\begin{cases}
        \text{若$\delta_2-\delta_3=-\frac{\pi}{2}$，称所构建的坐标系为右手系} \\
        \text{若$\delta_2-\delta_3=\frac{\pi}{2}$，称所构建的坐标系为左手系}
    \end{cases}\]
    基于右手系，令$\delta_2 \equiv 0$，则$\delta_3 = \frac{\pi}{2}$.
    
    故而$\hat{S}_z$表象下：
    $$\hat{S}_x = \frac{\hbar}{2} \begin{bmatrix} 0 & 1 \\ 1 & 0 \end{bmatrix} \qquad \hat{S}_y = \frac{\hbar}{2} \begin{bmatrix} 0 & -i \\ i & 0 \end{bmatrix} \qquad \hat{S}_z = \frac{\hbar}{2} \begin{bmatrix} 1 & 0 \\ 0 & -1 \end{bmatrix}$$
\end{derivation}
\begin{extension}{$\hat{S}_x$和$\hat{S}_z$的对易子$[\hat{S}_x,\hat{S}_z]$}
    \begin{align*}
        [\hat{S}_x, \hat{S}_z] &= \hat{S}_x \hat{S}_z - \hat{S}_z \hat{S}_x \\
        &= \frac{\hbar^2}{4} \begin{bmatrix} 0 & 1 \\ 1 & 0 \end{bmatrix} \begin{bmatrix} 1 & 0 \\ 0 & -1 \end{bmatrix} - \frac{\hbar^2}{4} \begin{bmatrix} 1 & 0 \\ 0 & -1 \end{bmatrix} \begin{bmatrix} 0 & 1 \\ 1 & 0 \end{bmatrix} \\
        &= \frac{\hbar^2}{4} \left( \begin{bmatrix} 0 & -1 \\ 1 & 0 \end{bmatrix} - \begin{bmatrix} 0 & 1 \\ -1 & 0 \end{bmatrix} \right) \\
        &= \frac{\hbar^2}{4} \begin{bmatrix} 0 & -2 \\ 2 & 0 \end{bmatrix} \\
        &= -i\hbar \cdot \frac{\hbar}{2} \begin{bmatrix} 0 & -i \\ i & 0 \end{bmatrix} \\
        &= -i\hbar \hat{S}_y
    \end{align*}
    经过不同坐标的计算，不难发现如下规律：
    $$[\hat{S}_\alpha,\hat{S}_\beta] = \varepsilon_{\alpha\beta\gamma} i\hbar\hat{S}_\gamma$$
    其中，$\varepsilon_{\alpha\beta\gamma}$称为\textbf{Levi-civita记号}，逆序数为偶时等于$1$，为奇时等于$-1$.
    
    还可以定义$\hat{S}^2 = \hat{S}_x^2 + \hat{S}_y^2 + \hat{S}_z^2 = \frac{3}{4}\hbar^2\hat{I}$，显然$\hat{S}^2$是个常数算符，而常数与任何算符的对易子均为$0$，有$[\hat{S}^2,\hat{S}_\alpha]=0$.
\end{extension}


\section{(不)相容力学量和对易关系}
\begin{definition}
    \begin{itemize}
        \item 称一个力学量$\hat{A}$的本征态对于另一个力学量$\hat{B}$而言是叠加态(从而$\hat{A},\hat{B}$无法同时取确定的值)的力学量为\textbf{不相容力学量}.
        \item 称存在一组共同的本征态力学量$\hat{A}$和力学量$\hat{B}$为\textbf{相容力学量}.
    \end{itemize}
\end{definition}

算符对应测量，若$\hat{A},\hat{B}$可同时取确定的值，则先测$\hat{A}$还是$\hat{B}$无区别.

故而有$\hat{A}\hat{B}=\hat{A}\hat{B},[\hat{A},\hat{B}]=0$.

\begin{proposition}
    当且仅当$[\hat{A},\hat{B}]=0$，力学量$\hat{A}$和力学量$\hat{B}$存在一组共同的本征态.
\end{proposition}
\begin{proof}
    \textbf{充分性：}

    设二者有共同本征态$\{\ket{n}\}$，即$\begin{cases} \hat{A}\ket{n} = a_n\ket{n} \\ \hat{B}\ket{n} = b_n\ket{n} \end{cases}$

    算符只有作用在一个态上才有意义：
    \begin{align*}
        \comm{\hat{A}}{\hat{B}}\ket{\psi} &= (\hat{A}\hat{B} - \hat{B}\hat{A})\hat{I}\ket{\psi} \\
        &= \sum_{n} (\hat{A}\hat{B} - \hat{B}\hat{A})\ket{n}\bra{n}\ket{\psi} \\
        &= \sum_{n} (\hat{A}b_n - \hat{B}a_n)\ket{n}\bra{n}\ket{\psi} \\
        &= \sum_{n} (a_n b_n - b_n a_n)\ket{n}\bra{n}\ket{\psi} \\
        &= 0 \quad \Rightarrow \comm{\hat{A}}{\hat{B}} = 0
    \end{align*}
    \rule{\textwidth}{0.4pt}
    \textbf{必要性：}

    设$\ket{n}$是$\hat{A}$的本征态：
    \begin{align*}
        0 &= \comm{\hat{A}}{\hat{B}}\ket{n} \\
        &= (\hat{A}\hat{B} - \hat{B}\hat{A})\ket{n} \\
        &= (\hat{A}\hat{B} - \hat{B}a_n)\ket{n} \quad \Rightarrow \hat{A}\hat{B}\ket{n} = a_n \hat{B}\ket{n}
    \end{align*}
    则$\hat{B}\ket{n}$是$\hat{A}$的本征矢，对应本征值$a_n$.
    \begin{itemize}
        \item \textbf{若$a_n$非简并：} \\
        于是$\hat{B}\ket{n}$与$\ket{n}$是同一个态.\\
        则$\hat{B}\ket{n}=b_n\ket{n}$(本征矢的长度没有意义). \\
        因此$\ket{n}$也是$\hat{B}$的本征态.
        \item \textbf{若$a_n$存在简并：} \\
        表明$a_n$对应的多个本征值构成了一个本征子空间，且经过$\hat{B}$作用的$\ket{n}$(即$\hat{B}\ket{n}$)对应本征值$a_n$. \\
        也就是说，$\hat{B}\ket{n}$也在$a_n$对应的本征子空间中. \\
        将$\hat{B}$在该本征子空间中的表象对角化，根据定理：
        \begin{center}
            \textit{算符在用其自身的一组正交归一本征矢构成的基下表示为一个对角矩阵}
        \end{center}
        对角化构造了该本征子空间中的一组新的基，且这组基是由$\hat{B}$的正交归一本征矢构成的. \\
        这个新基必然可以在$\ket{n}$所在的原基上被线性表出，而$a_n$所对应的本征矢的线性组合仍是对应$a_n$的本征矢. \\
        即新基也是对应本征值$a_n$的$\hat{A}$的本征矢，又是$\hat{B}$对应另一本征值$b_n$的本征矢. \\
        所以$\hat{A},\hat{B}$有一组共同的本征态.
    \end{itemize}
\end{proof}

\subsection{基本对易关系}
计算$\comm{x}{\hat{P}_X}$：

将$\comm{x}{\hat{P}_X}$作用在任意波函数上：
\begin{align*}
    \comm{x}{\hat{p}_x} \psi(x) &= (x \hat{p}_x - \hat{p}_x x) \psi(x) \\
    &= -i\hbar \left\{ x \frac{d}{dx} \psi(x) - \frac{d}{dx} [x \psi(x)] \right\} \\
    &= -i\hbar \left\{ x \frac{d}{dx} \psi(x) - \left[ \psi(x) + x \frac{d}{dx} \psi(x) \right] \right\} \\
    &= i\hbar \psi(x) \quad \Rightarrow \comm{x}{\hat{p}_x} = i\hbar \quad \text{基本对易关系}
\end{align*}

\subsection{对易式基本性质}
\begin{itemize}
    \item $\comm{\hat{A}}{\hat{B}} = -\comm{\hat{B}}{\hat{A}}$
    \item $\comm{\hat{A}}{f(\hat{A})} = 0$
    \item $\comm{\hat{A}}{\text{const}} = 0$
    \item $\comm{\hat{A}}{\hat{B} + \hat{C}} = \comm{\hat{A}}{\hat{B}} + \comm{\hat{A}}{\hat{C}}$
    \item $\comm{\hat{A}}{\hat{B}\hat{C}} = \hat{B}\comm{\hat{A}}{\hat{C}} + \comm{\hat{A}}{\hat{B}}\hat{C}$
    \item $\comm{\hat{A}\hat{B}}{\hat{C}} = \hat{A}\comm{\hat{B}}{\hat{C}} + \comm{\hat{A}}{\hat{C}}\hat{B}$
    \item $\comm{\hat{A}}{\comm{\hat{B}}{\hat{C}}} + \comm{\hat{B}}{\comm{\hat{C}}{\hat{A}}} + \comm{\hat{C}}{\comm{\hat{A}}{\hat{B}}} = 0$ \quad 雅可比恒等式
\end{itemize}
\begin{extension}{}
    \begin{itemize}
        \item \textbf{计算$\comm{x}{-\frac{\hbar^2}{2\mu} \hat{p}^2 + V(x)}$}
        \begin{align*}
            \comm{x}{-\frac{\hbar^2}{2\mu} \hat{p}^2 + V(x)} &= -\frac{\hbar^2}{2\mu} \comm{x}{\hat{p}^2} + \comm{x}{V(x)} \\
            &= -\frac{\hbar^2}{2\mu} (\hat{p} \comm{x}{\hat{p}} + \comm{x}{\hat{p}} \hat{p}) \\
            &= -\frac{i\hbar^3}{\mu} \hat{p}
        \end{align*}
        \item \textbf{计算$\comm{\hat{p}}{f(x)}$}
        \begin{align*}
            \comm{\hat{p}}{f(x)} \psi(x) &= -i\hbar \left[ \pdv{x} f(x) - f(x) \pdv{x} \right] \psi(x) \\
            &= -i\hbar \left[ \pdv{x} (f(x) \psi(x)) - f(x) \pdv{x} \psi(x) \right] \\
            &= \left[ -i\hbar \pdv{x} f(x) \right] \psi(x) \\
            \Rightarrow \comm{\hat{p}}{f(x)} &= -i\hbar \pdv{x} f(x) = \comm{\hat{p}}{x} \pdv{x} f(x)
        \end{align*}
        \item \textbf{计算$\comm{x}{f(\hat{p})}$} \\
        \[\comm{x}{f(\hat{p})} = \comm{x}{\hat{p}} \frac{\partial}{\partial \hat{p}} f(\hat{p})\]
    \end{itemize}
\end{extension}


\section{力学量完全集(CSCO)}



\section{Heisenberg不确定性原理和Ehrenfest定理}



\section{表象变换}



\section{动量表象波函数}



\section{动量表象位置算符}



\section{动量表象下的能量本征方程}



\section{动量与空间平移算符}